\section{Data and reproducibility}
\label{app:data}

Using LLM-assisted web search, we collected the \PractitionerAccountCount{}
public first-person accounts used in this paper.  Each row retains the public
source URL, bibliography key, observation used here, a qualification note, and
the categorical fields used for the descriptive counts.  The checked-in source
notes render these rows as per-account records so each observation can be
checked against its public source.  The collection is not necessarily
representative, so we make no prevalence estimates from it.

The account table is accompanied by a screening log.  The log records the
sources in the preexisting snapshot and candidates reviewed from 8 September
2026 onward, with their disposition and a reason.  It does not reconstruct every
candidate encountered in earlier searches and is not an exhaustive search
record.

\texttt{make sources} validates the account and screening schemas, cross-table
identifiers, citation keys, and URLs and regenerates the paper macros.  Figure~\ref{fig:lean-activity}
reproduces the Papers With Lean public monthly series for January 2025 through
August 2026 from a 6 September 2026 snapshot.  Grouping the captured corpus by
its \texttt{published} month reproduced the public chart; the partial September
count is excluded from the figure.

\ifshowartifacturl
Project materials: \href{\artifacturl}{public GitHub repository}.\\
Lean helper tooling, including \texttt{leanq}: \href{\leanqurl}{aiq-lean-formalization-tools}.
\else
Identifying repository URLs are suppressed in this double-blind build and are
included in the supplementary artifact.
\fi

The build records SHA-256 digests for manuscript inputs and generated evidence
in \texttt{generated/materials\_manifest.json}.

The 29-result register counts the four unnumbered Section~2 headline theorems
and every named theorem, proposition, lemma, and corollary established in the
source.  Definitions, proof steps, examples, immediate consequences, and open
or deferred claims are not counted separately.  Each row records the source
location and scope, the Lean declaration used as evidence, and the review state.
Proposition~4.4 remains a counted result: its formal treatment certifies a
counterexample to the printed claim and proves a repaired theorem.

\section{Davis--Kahan background and Formalization 1}
\label{app:formalization1}

Davis and Kahan compare a candidate subspace with an exact invariant
subspace of a self-adjoint operator \citep{davis1970rotation}.  In finite
dimensions, let $A=A^*$ and let $E_0$ have orthonormal columns spanning the
candidate subspace $V$.  The map $E_0$ takes coordinates in $V$ to the ambient
space.  Applying $A$ and then returning to those coordinates gives the
compression $A_0=E_0^*AE_0$.  With $P_V=E_0E_0^*$, its residual is
\[
  R=AE_0-E_0A_0=(I-P_V)AE_0.
\]
Thus $R$ is the component of $AE_0$ outside $V$, and $R=0$ precisely when $V$
is invariant under $A$.  The historical signature allows any bounded operator
\texttt{M} on $V$, not just this compression; its residual is
$A\circ\iota_V-\iota_V\circ M$, where $\iota_V:V\hookrightarrow E$ is inclusion.
Choosing $M$ to be the compression recovers the preceding interpretation.

For the exact subspace $U$, the assumption \texttt{Reduces A U} says that both
$U$ and $U^\perp$ are invariant.  In Formalization~1, the spectrum of $A|_U$
lies in $[a,b]$, while the spectrum of $A|_{U^\perp}$ lies outside
$(a-\delta,b+\delta)$ for $\delta>0$.  This gap separates the two exact spectral
parts.  It is an assumption on the operator being compared, rather than a
consequence of the residual being small.

For two subspaces of the same finite dimension with orthonormal coordinate
maps $E_0$ and $F_0$, their principal angles $\theta_j\in[0,\pi/2]$ satisfy
$\cos\theta_j=\sigma_j(F_0^*E_0)$.  Equivalently, they pair directions in the two
subspaces by their overlap.  The directed sine measures the part of a candidate
direction outside the exact subspace; its matrix representative is
$(I-P_U)E_0$.  The doubled-angle quantity applies $\sin(2\theta_j)$ to the
angles, with the candidate-to-exact orientation fixed by the source.  The bound
\[
  \delta\,N(\sin 2\Theta_0)\le 2N(R)
\]
therefore controls the doubled-angle quantity by the residual relative to the
gap.  A small doubled sine alone does not imply a small angle:
$\sin(2\theta)$ vanishes at both $0$ and $\pi/2$.  Recovering a small-angle bound
requires an appropriate angle branch, such as $0\le\theta\le\pi/4$.

Formalization~1 states the estimate for complex Hilbert spaces.  Its ambient
Hilbert-space assumptions are omitted from the displayed code.  The type
\texttt{ContinuousLinearMap (RingHom.id Complex) E E} restricts $A$ to a bounded
complex-linear operator.  The identity scalar homomorphism specifies complex
linearity.  The following table connects the remaining structures to the
mathematics.

\begin{table}[H]
\centering
\small
\caption{Reader's guide to the principal Lean structures in Formalization~1.}
\label{tab:formalization1-glossary}
\begin{tabularx}{\linewidth}{@{}>{\raggedright\arraybackslash}p{0.36\linewidth}>{\raggedright\arraybackslash}X@{}}
\toprule
Lean expression & Mathematical role \\
\midrule
\texttt{ContinuousLinearMap (RingHom.id Complex) E E}
  & A bounded continuous complex-linear operator $E\to E$.  In Formalization~1
    this is the ambient self-adjoint operator $A$; restricting it to this type
    is part of the historical scope mismatch. \\
\texttt{IsSelfAdjoint A}
  & $A=A^*$, the self-adjointness assumption. \\
\texttt{Submodule Complex E}
  & A complex linear subspace of the ambient Hilbert space.  Here $U$ is the
    exact reducing subspace and $V$ is the candidate subspace. \\
\texttt{U.HasOrthogonalProjection}
  & The orthogonal projection onto $U$ is available; likewise for $V$.  This
    supports $U^\perp$, compressions, and the angle constructions. \\
\texttt{Reduces A U}
  & Both $U$ and $U^\perp$ are invariant under $A$. \\
\texttt{compressOperator U A}
  & The orthogonal compression $P_U A|_U$; for a reducing subspace this is the
    restriction of $A$ to $U$. \\
\texttt{spectrum Real (...)}
  & The real spectrum of the self-adjoint compression, used to state the
    interval and gap hypotheses. \\
\texttt{V.subtypeL}
  & The bounded inclusion $V\hookrightarrow E$, corresponding to the source's
    coordinate map $E_0$. \\
\texttt{M : ContinuousLinearMap (RingHom.id Complex) V V}
  & An arbitrary bounded operator on the candidate subspace.  Taking it to
    be the compression gives the introductory choice of $A_0$. \\
\texttt{residual A V.subtypeL M}
  & $A\circ\iota_V-\iota_V\circ M$.  The inclusion plays the role of $E_0$
    and $M$ the role of the candidate-space operator $A_0$. \\
\texttt{SymmetricNormingFunction}
  & A normalized, dimension-coherent symmetric gauge on singular values used to
    represent the source's unitarily invariant norm $N$. \\
\texttt{N.Mem T}
  & $T$ lies in the canonical operator ideal on which the norm determined by
    $N$ is finite. \\
\texttt{N.gauge T}
  & The real-valued gauge used as $N(T)$ in the inequality.  The accompanying
    membership hypothesis places $T$ in its normed ideal. \\
\texttt{sinTwoThetaIdealBlock U V}
  & A projection-product representative of the doubled-angle quantity.
    Its approximation numbers agree with those of the directed double-angle
    sine, which identifies the norms in the estimate. \\
\bottomrule
\end{tabularx}
\end{table}

The conclusion asserts that the doubled-angle block belongs to the same ideal
as the residual and then bounds its gauge.  These norm and spectral conditions
are part of the statement.  The historical scope mismatch is visible in the
fixed complex field and bounded type of $A$.  They exclude the source's real
and unbounded settings, even when all the displayed hypotheses hold.

\section{Lean publication activity}

\begin{figure}[H]
  \centering
  \input{generated/lean_activity_timeline_tikz.tex}
  \caption{Monthly publication counts reproduced from the Papers With Lean
  public statistics series, January 2025 through August 2026
  \citep{paperswithlean2026stats}.  Counts are grouped by the captured data's
  \texttt{published} month, and only complete months are shown.  The fraction
  that used AI is unknown.}
  \label{fig:lean-activity}
\end{figure}

\section{Project chronology}

\begin{table}[H]
\centering
\small
\caption{Selected public Git events used to anchor the chronology.  The CSV
stores the full commit identifiers and evidence paths; the build validates both.}
\label{tab:review-timeline}
\begin{tabularx}{\linewidth}{@{}ll>{\raggedright\arraybackslash}X@{}}
\toprule
Date & commit & event \\
\midrule
\input{generated/review_timeline_table.tex}
\end{tabularx}
\end{table}

\section{Semantic-review page}

\begin{figure}[H]
  \centering
  \includegraphics[width=\linewidth,trim=0 250 0 0,clip]{figures/semantic-alignment-sine-theta-row.png}
  \caption{A review interface developed late in the project.  The source passage
  and inherited context appear at left, the claimed source-to-Lean
  correspondence in the middle, and compiler-derived Lean evidence at right.
  Most reviews reported in the main paper used LLM source/signature comparison;
  the interface was added later.}
  \label{fig:dashboard}
\end{figure}

The interface consolidates source passages, correspondence judgments, and
compiler-derived signatures previously inspected through files and Lean
queries.  Understanding a signature can also require the definitions of its
structures, predicates, coercions, and typeclass assumptions.  Lean Compass's
semantic dependency view suggests how these definitions could be exposed while
omitting proof-only dependencies \citep{yanahama2026leanatlas}.

\section{Resource transparency}
\label{app:resources}

To estimate the cost of the formalization effort, we maintained a resource
accounting record.  The compact snapshot below has an accounting cutoff of
13 August 2026.  Retained model turns and token quantities are incomplete
measured lower bounds; price and serving-energy rows are estimates derived from
those observations.  The energy estimate excludes training, embodied hardware,
client devices, and network energy.

\begin{table}[H]
\centering
\scriptsize
\caption{Observed lower bounds and derived estimates at the 13 August 2026
accounting cutoff.  These are not project totals.}
\begin{tabularx}{\linewidth}{@{}p{0.28\linewidth}p{0.17\linewidth}>{\raggedright\arraybackslash}X@{}}
\toprule
Quantity & Lower bound / estimate & Basis / interpretation \\
\midrule
\input{generated/resource_snapshot_table.tex}
\end{tabularx}
\end{table}


\begin{table}[H]
\centering
\scriptsize
\caption{AI systems and model families used during the formalization.
Token telemetry is incomplete.}
\label{tab:formalization-systems}
\begin{tabularx}{\linewidth}{@{}p{0.20\linewidth}p{0.27\linewidth}>{\raggedright\arraybackslash}X@{}}
\toprule
System & Models used & Typical use \\
\midrule
\input{generated/model_systems_table.tex}
\end{tabularx}
\end{table}

\begin{table}[H]
\centering
\scriptsize
\caption{Recovered model-resolved token measurements at the 13 August 2026
accounting cutoff.  These are measured lower bounds for the five labels with
recovered telemetry; models in \cref{tab:formalization-systems} that do not
appear here were still used.}
\label{tab:model-tokens-appendix}
\resizebox{\linewidth}{!}{% Generated by scripts/build_accounting.py; do not edit by hand.
\begin{tabular}{lrrrr}
\toprule
Model & \multicolumn{4}{c}{Retained measured tokens} \\
 & Input & Cache write & Cache read & Output \\
\midrule
claude-opus-5 & 140,093 & 116,791,540 & 17,395,271,984 & 34,338,937 \\
claude-opus-4-8 & 693,366 & 58,990,139 & 3,744,966,700 & 15,967,324 \\
claude-fable-5 & 380,759 & 27,789,153 & 1,512,917,593 & 7,985,843 \\
gpt-5.6-sol & 6,502,664 & 0 & 343,618,944 & 635,601 \\
gpt-5.6-terra & 3,798,915 & 0 & 382,951,936 & 563,722 \\
claude-sonnet-5 & 78 & 110,211 & 2,256,550 & 26,074 \\
\bottomrule
\end{tabular}
}
\end{table}

Historical model labels are retained verbatim, with exact recovered telemetry used
when it is available.
